\section*{Bab \ref{ch:g10:habitable-planet} --- Planet yang Layak Huni}

\begin{solution}{exo:g10:habitable-planet:1}
Suhu ketika pemancar sempurna akan memancarkan ke angkasa daya yang tepat sama
dengan daya yang diserap planet itu dari bintangnya. Permukaan Bumi lebih
hangat (\qty{288}{K} dibandingkan dengan sekitar \qty{255}{K}) karena gas
rumah kaca menyerap sebagian inframerah yang dipancarkannya dan mengirim
sebagian kembali ke bawah.
\end{solution}

\begin{solution}{exo:g10:habitable-planet:2}
Venus: terukur \qty{737}{K}, sedangkan \omterm{def:g10:habitable-planet:equilibrium-temperature}{suhu kesetimbangannya} sekitar
\qty{227}{K}, jadi pemanasannya sekitar \qty{510}{K}.
\end{solution}

\begin{solution}{exo:g10:habitable-planet:3}
$S = 1361/1.52^2 \approx \qty{589}{W/m^2}$ (lembar fakta memberikan
\qty{586}{W/m^2}, dengan jarak rata-rata yang tepat).
\end{solution}

\begin{solution}{exo:g10:habitable-planet:4}
\qty{1}{bar}, \qty{330}{K}: cair (di bawah titik didih, sekitar
\qty{373}{K}). \qty{0.001}{bar}, \qty{250}{K}: es (tekanannya berada di atas
kurva es--uap, yang pada \qty{250}{K} terletak sekitar \qty{0.0008}{bar}).
\qty{100}{bar}, \qty{700}{K}: uap --- suhunya di atas suhu kritis, jadi air
cair tidak mungkin ada.
\end{solution}

\begin{solution}{exo:g10:habitable-planet:5}
Planet yang mengorbit bintang selain Matahari. Transit: planet melintas di
depan bintangnya dan sedikit meredupkannya pada setiap orbit. Goyangan:
tarikan planet menggerakkan bintang bolak-balik di sepanjang garis pandang
kita, sehingga warna cahayanya bergeser secara berkala.
\end{solution}

\begin{solution}{exo:g10:habitable-planet:6}
$T = (1361 \times 0.706/(4 \times 5.67\times10^{-8}))^{1/4} \approx
\qty{255}{K} \approx \qty{-18}{\degreeCelsius}$. Dengan $A = 0$:
$\approx \qty{278}{K} \approx \qty{5}{\degreeCelsius}$.
\end{solution}

\begin{solution}{exo:g10:habitable-planet:7}
$\sqrt 4 = 2$: dari $2 \times 0.99 = \qty{1.98}{au}$ sampai $2 \times 1.70 =
\qty{3.40}{au}$. Pada \qty{1}{au} planet itu akan menerima empat kali
iradiansi Bumi: jauh di sebelah dalam tepi dalam, terlalu panas.
\end{solution}

\begin{solution}{exo:g10:habitable-planet:8}
$S = 1361/0.723^2 \approx \qty{2604}{W/m^2}$;
$T_{\mathrm{eq}} = (2604 \times 0.23/(4\sigma))^{1/4} \approx \qty{227}{K}$.
Permukaannya sekitar \qty{510}{K} lebih panas: atmosfer karbon dioksidanya
yang tebal (\qty{92}{bar}) menghasilkan efek rumah kaca yang sangat besar.
\end{solution}

\begin{solution}{exo:g10:habitable-planet:9}
Pada \qty{6.36}{mbar} air hanya berwujud cair di antara sekitar
\qty{0}{\degreeCelsius} dan \qty{0.5}{\degreeCelsius}. Pada
\qty{10}{\degreeCelsius} tekanannya berada di bawah kurva didih: air mendidih
dan menguap; pada malam hari air membeku.
\end{solution}

\begin{solution}{exo:g10:habitable-planet:10}
$11.2/2.4 \approx 4.7$. Molekul gas harus bergerak hampir lima kali lebih
cepat untuk lepas dari Bumi daripada dari Bulan; selama miliaran tahun
gravitasi Bulan yang lemah membiarkan gas-gasnya lepas, sedangkan gravitasi
Bumi mempertahankan sebagian besar gasnya.
\end{solution}

\begin{solution}{exo:g10:habitable-planet:11}
$T = (1361 \times 0.50/(4\sigma))^{1/4} \approx \qty{234}{K}$, sekitar
\qty{-39}{\degreeCelsius}. Es akan memantulkan begitu banyak cahaya Matahari
sehingga planet tetap cukup dingin untuk mempertahankan es itu: sebuah
lingkaran yang saling memperkuat.
\end{solution}

\begin{solution}{exo:g10:habitable-planet:12}
$T_{\mathrm{eq}} = (S_0(1-A)/(4\sigma d^2))^{1/4} = (S_0/(4\sigma))^{1/4}
(1-A)^{1/4} (d^2)^{-1/4}$; $(1361/(4\sigma))^{1/4} \approx \qty{278.3}{K}$ dan
$(d^2)^{1/4} = \sqrt d$. Empat kali lebih jauh: $T$ dibagi $\sqrt 4 = 2$.
\end{solution}

\begin{solution}{exo:g10:habitable-planet:13}
$\sqrt d = 278.3 \times 0.706^{1/4}/273 \approx 0.934$, jadi $d \approx
\qty{0.873}{au}$, di sebelah dalam tepi dalam (\qty{0.99}{au}). \omterm{def:g10:habitable-planet:equilibrium-temperature}{Suhu
kesetimbangan} mengabaikan efek rumah kaca: planet sungguhan yang beratmosfer
lebih hangat daripada $T_{\mathrm{eq}}$, jadi planet itu sudah kepanasan dan
kehilangan airnya pada jarak yang lebih jauh, sekitar \qty{0.99}{au}.
Tepi-tepi zona berasal dari model seluruh atmosfer, bukan dari
$T_{\mathrm{eq}}$ saja.
\end{solution}

\begin{solution}{exo:g10:habitable-planet:14}
Di antara titik leleh, sekitar \qty{0}{\degreeCelsius}, dan titik didih,
sekitar \qty{7}{\degreeCelsius}. Pada \qty{15}{\degreeCelsius} danau itu
akan mendidih.
\end{solution}

\begin{solution}{exo:g10:habitable-planet:15}
$\sqrt{0.5} \approx 0.707$: dari $0.707 \times 0.99 \approx \qty{0.70}{au}$
sampai $0.707 \times 1.70 \approx \qty{1.20}{au}$. Planet itu tepat berada di
tepi luar. Tepi-tepi zona adalah hasil model dan penskalaannya hanya
pendekatan, jadi mustahil dikatakan apakah planet itu di dalam atau di luar.
\end{solution}

\begin{solution}{pb:g10:habitable-planet:1}
\textbf{1.} $0.00151/0.0485^2 \approx 0.642$.

\textbf{2.} $0.642 \times 1361 \approx \qty{874}{W/m^2}$.

\textbf{3.} Di antara keduanya: sekitar tiga kali lebih kecil daripada di
Venus, satu setengah kali lebih besar daripada di Mars.

\textbf{4.} $365.25/11.18 \approx 32.7$.

\textbf{5.} $(874/(4 \times 5.67\times10^{-8}))^{1/4} \approx \qty{249}{K}$.

\textbf{6.} $(874 \times 0.23/(4\sigma))^{1/4} \approx \qty{173}{K}$.

\textbf{7.} Tidak ada: \qty{249}{K} pun masih di bawah \qty{273}{K}. Selubung
awan atau es yang terang menurunkan \omterm{def:g10:habitable-planet:equilibrium-temperature}{suhu kesetimbangan} lebih dari
\qty{75}{K}: apa yang menutupi sebuah planet sama pentingnya dengan jaraknya.

\textbf{8.} $\sqrt{0.00151} \approx 0.0389$.

\textbf{9.} $0.99 \times 0.0389 \approx \qty{0.0385}{au}$.

\textbf{10.} $1.70 \times 0.0389 \approx \qty{0.066}{au}$.

\textbf{11.} Ya: $0.0385 < 0.0485 < 0.066$.

\textbf{12.} $d_\odot = 0.0485/0.0389 \approx \qty{1.25}{au}$: di paruh dalam
\omterm{def:g10:habitable-planet:habitable-zone}{zona layak huni} Matahari, di antara Bumi dan Mars.

\textbf{13.} $1.055^{1/3} \approx 1.018$.

\textbf{14.} $11.2 \times \sqrt{1.055/1.018} \approx \qty{11.4}{km/s}$, kira-kira
sama dengan laju lepas Bumi: gravitasinya dapat menahan atmosfer (apakah
planet itu memang memilikinya bergantung pada sejarahnya).

\textbf{15.} Goyangan bintang hanya mengungkapkan bagian geraknya yang
searah garis pandang kita; karena kemiringan orbitnya tidak diketahui, massa
yang diperoleh sama dengan massa sebenarnya hanya jika orbit itu terlihat dari
samping, dan selebihnya merupakan nilai minimum.

\textbf{16.} Tekanan titik tripel, \qty{611.657}{Pa}, sekitar
\qty{0.006}{bar}.

\textbf{17.} $(874 \times 0.706/(4\sigma))^{1/4} \approx \qty{228}{K}$.

\textbf{18.} $288 - 228 = \qty{60}{K}$: Proxima~b memerlukan pemanasan rumah
kaca \textbf{sekitar \qty{60}{K}}, hampir dua kali \qty{33}{K} milik Bumi.
\end{solution}
